
\begin{abstract}

水泥窑尾电除尘器的协同优化控制关乎超低排放达标与运行能耗的平衡，是当前水泥行业绿色转型中的关键工程问题。本研究以某$5000$吨/日新型干法生产线窑尾四电场电除尘器为对象，利用连续$7$天分钟级运行数据，体系化地考察了入口工况与操作参数对出口排放及电耗的定量影响、超低排放约束下的分工况寻优策略、调控参数优先级判别以及排放标准收紧的电耗代价评估。

\textbf{问题一}围绕参数关系建模展开。数据分析揭示$C_{out}$标准差仅$0.17$且数据最大值被限制为$50\;\text{mg/Nm}^3$，系传感器量程限幅所致，据此确立机理主导的建模路线。以Deutsch方程为骨架构建多级串联效率模型，引入振打周期双向偏离指数衰减修正；电耗模型则基于电场能量$\propto U^2$与振打功耗$\propto 1/T$的物理关系建立。Deutsch参数经对数空间残差L-BFGS-B多起点（$50$次）拟合，电耗系数经约束最小二乘辨识，随机森林特征重要性用于交叉验证。所得扩展电耗模型$R^2=0.9979$，Deutsch参数$kA_0=287.25$；各特征重要性均匀且$p>0.05$，从数据层面印证了限幅下纯数据驱动外推的不可靠性。

\textbf{问题二}解决工况划分与约束寻优。以$(C_{in},T_{in},Q)$为特征在标准化空间执行K-Means聚类，轮廓系数择定$K=6$；各工况下建立$\min P_{total}$为目标、$C_{out}\leq 10\;\text{mg/Nm}^3$为约束的8维非线性规划，SLSQP多起点（$10$个，种子$42$）求解，不收敛则回退至差分进化全局搜索。6个工况全部成功求解且排放达标，最优电耗$1831\sim1993\;\text{kW}$，高浓度工况电耗明显偏高，与物理预期吻合。

\textbf{问题三}聚焦灵敏度与优先级判别。取最高浓度（$C_{in}=44.34$）与最低浓度（$C_{in}=26.01$）工况对比，中心差分法计算$C_{out}$与$P_{total}$对各参数的数值偏导，步长减半校验一致性，以无量纲弹性性价比$|E^C/E^P|$判定优先级。纯电耗模型下振打边际性价比$29.37$约为电压$12.95$的2.3倍（根因：$\partial P/\partial T$远小于$\partial P/\partial U$，分母小而非振打降排放能力强），引入振打磨损隐性成本修正因子$\gamma=2.5$后反转为电压优先（临界$\gamma^*=2.27$），推荐\textbf{"电压为主、振打微调"的分层调控策略}。

\textbf{问题四}评估排放收紧代价。将约束收紧至$C_{out}\leq 5\;\text{mg/Nm}^3$复用优化框架重新求解，逐工况计算电耗增幅$\Delta P\%$并校验物理可行性。整体平均增幅$4.95\%$，各工况$4.80\%\sim5.15\%$，全部工况参数落在可行域内。

本研究贯通了机理建模、工况优化、灵敏度分析及排放收紧评估的研究链条。须坦诚说明，$C_{out}$限幅致历史数据集中于$50\;\text{mg/Nm}^3$附近，由此外推至$10$乃至$5\;\text{mg/Nm}^3$的定量结论承载较大不确定性，后续需经主动实验验证。研究成果可为水泥电除尘器超低排放协同优化提供趋势性参考与工程决策支撑。\footnote{AI辅助润色}

\keywords{电除尘器；Deutsch方程；约束优化；K-Means工况划分；灵敏度分析}

\end{abstract}
